\section{Tighter Regret Bounds with Duality and Local Norms}
\begin{frame}{}
    We now shift away from using the FoReL framework to study the regret bounds of various OMD algorithms and instead adopt a new lens to perhaps get tighter bounds and in a much simpler way. In this section, we introduce two other ways to analyse regret bounds.
    \begin{itemize}
        \item Duality
        \item Local Norms
    \end{itemize}
\end{frame}

\begin{frame}{A Primer on Fenchel Conjugacy}
We first give the definition of the Fenchel conjugate of a convex function $f$:
\begin{block}{Fenchel Conjugate of a convex function $f$}
    The Fenchel conjugate of the convex function $f$, denoted as $f^*$ is defined to be\begin{equation}\label{eqn:frenchel_conjugate}
        f^*(\theta) = \sup_{x \in \text{dom}(f)} \langle \theta, x \rangle -f(x)
    \end{equation}
\end{block}
Intuitively, we can think of the Fenchel conjugate of $f$ as an alternative representation of a function.\pause
\begin{equation*}
    \text{Set of points} : (x, f(x)) \to \text{Set of tangents} : (\theta, -f^*(\theta))
\end{equation*}\pause
The Fenchel conjugate $f^*$ is the function which takes in $\theta$ and outputs the \textbf{negative} $y$-intercept of the tangent line with the slope $\theta$, $f^*(\theta)$.
    
\end{frame}

\begin{frame}{Fenchel Conjugate Pairs examples}
\begin{table}[h]
\centering
\renewcommand{\arraystretch}{1.3}
\begin{tabular}{ccl}
\toprule
$f(x)$ & $f^\star(\theta)$ & Comments \\
\midrule
$\frac{1}{2}\|x\|_2^2$ & $\frac{1}{2}\|\theta\|_2^2$ & \\
$\frac{1}{2}\|x\|_q^2$ & $\frac{1}{2}\|\theta\|_p^2$ & where $\frac{1}{p} + \frac{1}{q} = 1$ \\
$\sum_i x[i]\log(x[i]) + I_{\Delta_d}(x)$ & $\log\big(\sum_i e^{\theta[i]}\big)$ & (normalized Entropy) \\
$\frac{1}{\eta} g(x)$ & $\frac{1}{\eta} g^\star(\eta\theta)$ & where $\eta > 0$ \\
$g(x) + \langle x, x' \rangle$ & $g^\star(\theta - x')$ & \\
\bottomrule
\end{tabular}
\caption{Some useful Fenchel conjugate pairs.}
\label{tab:fenchel-pairs}
\end{table}
\end{frame}

\begin{frame}{Strong-Smoothness}
    Recall the definition of $\sigma$-strong convexity. \pause
    \begin{block}{$\sigma$-strong convexity}
        A function $f$ that is $\sigma$-strongly convex over a set $S$ with respect to norm $||\cdot ||$ if for any $x' \in S$, we have that \begin{equation*}
        \forall z \in \partial f(x'),\;\forall x \in S, \; f(x) \geq f(x') + \langle z, x - x' \rangle + \frac{\sigma}{2} || x - x' ||^2
    \end{equation*}
    \end{block}
    \pause
    A related property of strong-convexity is what we call \textit{strong-smoothness} and it is related to the Bregman divergence
    \begin{block}{$\sigma$-Strong-smoothness}
        A function $h$ is $\sigma$-strong smooth with respect to a norm $||\cdot||$ if it is differentiable and for all $x, x'$, we have that
        \begin{equation}
            D_h(x||x') \leq \frac{\sigma}{2}||x - x'||^2
        \end{equation}
    \end{block}
\end{frame}

\begin{frame}{Strong-convexity-Strong-smoothness duality}
    \begin{block}{Strong-convexity-smoothness Duality Lemma}
        Assume that $h$ is a closed and convex function. $h$ is $\sigma$-strongly convex with respect to the norm $||\cdot||$ if and only if the Fenchel conjugate $h^*$ is $\frac{1}{\sigma}$-strongly smooth with respect to the dual norm $||\cdot||_*$.
    \end{block}
\end{frame}